The open-source LLM guardrail ecosystem in 2026 looks mature on the
surface. Content classifiers, dialogue DSLs, decorator-style
validator hubs, red-team frameworks: most of the squares on the
matrix are filled. Three squares are not. First, no production
library implements capability-based information-flow control around
tool-using agents, despite the CaMeL research line that drives
attack success rates to roughly zero by construction. Second, streaming guardrails are still wrong by
default: existing libraries either buffer the full response and
sacrifice time-to-first-token, or run guards in parallel and let
violative tokens leak before cancellation. Third, the round-trip
from PyRIT or Garak findings to deployed runtime guards is a
two-engineer-day workflow that Galileo and Lakera sell as a
commercial moat, and that has no open-source equivalent. We built
AIsafePy to fill these three gaps. The library has a small
dependency-light core and three focused modules.
\texttt{aisafepy.flow} ports CaMeL-style taint propagation to
OpenAI Agents SDK, LangGraph, LlamaIndex, Anthropic tools, and the
Model Context Protocol. \texttt{aisafepy.stream} is a three-tier
cascade with an explicit p95 latency budget and optional white-box
activation probes for self-hosted models. \texttt{aisafepy.adapt}
compiles red-team failures into four output modalities: distilled
classifiers, synthesized regexes, Cedar or OPA policies, and
contrastive steering vectors. We describe the design, key
implementation decisions, an evaluation roadmap against AgentDojo
and HarmBench plus a streaming-guard latency harness, and the
failure modes the library does not address.
